\section{Real scalp EEG}
\label{sec:real}

\begin{figure*}[t]
  \centering
  \includegraphics[width=\textwidth]{fig3_real_eeg.pdf}
  \caption{Real EEG against the benchmark. (a) Held-out recovery against the magnitude of the segment-wise modulation (standard deviation of the per-segment log-variance, in dB; log axis) for the source-specific schedule, and for TCL under one shared gain (dashed); chance band in grey; mean and bootstrap 95\% interval over 5 seeds. Shaded: the range of the nine BCI IV-2a subjects' observed modulation magnitude when the recording is cut into 20, 200 or 2000 segments (2~min, 12~s, 1.2~s per segment). (b) Four-class cue decoding accuracy from the per-trial log-variance of TCL components fit without labels on three runs and scored on the other three, against TCL's segment length; one point per subject, lines at the median, filled $M = 4$, hollow $M = 8$. Dotted: medians of PCL, which does not use segments, and of supervised CSP on all 22 channels; dashed: label-permutation 95th percentile. (c) Agreement between TCL fits on different runs, with the PCL and FastICA medians. Protocols in Appendix~\ref{app:real}.}
  \label{fig:real}
\end{figure*}

Two things can be done with the recording itself: check the conditions the theorems make checkable, and score the components against the one reference the experiment provides.

\paragraph{Real EEG is in the regime where the benchmark's conclusions hold.}
Two properties of the segment-wise modulation are observable without ground truth: how much of it follows one spatial pattern (the co-modulation fraction of Section~\ref{sec:benchmark}) and how large it is (the standard deviation of the per-segment log-variance).
On the nine BCI IV-2a subjects the co-modulation fraction lies in the left half of axis (d) of Figure~\ref{fig:axes}, in the broad band as in the motor band, away from the shared-gain regime (Table~\ref{tab:mod}); the rank condition is not refuted.
The magnitude depends on how the recording is segmented.
Figure~\ref{fig:real}a repeats the modulation cells of the benchmark at smaller magnitudes and marks where the subjects fall when their recordings are cut into 20, 200 or 2000 segments.
At the second-scale segmentation that TCL's own protocol uses, the magnitude of real EEG's modulation is the magnitude at which every conclusion of Section~\ref{sec:truth} holds, including the shared-gain failure; the benchmark cells use 20 segments of 0.4~s, so the segment length is comparable and the recordings supply many more segments.
At minute-scale segments the modulation is too weak for any segment-based method, which is a property of the segmentation rather than of the recording.
A TCL recovery reported on these recordings at its natural segment length would therefore pass every check the recording allows, and the failures that remain are the ones that cannot be seen in the recovery: drift, which preprocessing can remove but no score reveals, and nuisance processes with a spatial pattern of their own (Appendix~\ref{app:robust}).

\paragraph{The cue is a reference, and the components barely carry it.}
BCI IV-2a has an experimental manipulation: a randomised four-class motor-imagery cue with a known cortical locus.
We fit each unsupervised method on three runs of a subject without labels and score its components on the other three by how well their per-trial log-variance decodes the cue (Appendix~\ref{app:cue}); a label-permutation null and a supervised ceiling, common spatial patterns (CSP), bracket the result.
Figure~\ref{fig:real}b,c follow TCL across the same three segment lengths; the other methods do not use segments and enter as reference lines.
Finer segments help TCL, as panel (a) predicts: its decoding accuracy and its agreement between fits on different runs both rise.
At its natural segment length, however, TCL's components decode the cue no better than the components of methods that never use the segments (paired Wilcoxon tests over the nine subjects against PCA, FastICA and PCL, none significant; Appendix~\ref{app:cue}), well below CSP ($p = 0.004$), and at every segment length they remain the least reproducible of the unsupervised methods ($p = 0.004$ against each).
On real EEG the segment-based method, at its best setting, finds nothing about the only available reference that a linear baseline does not.
